\setcounter{ex}{0} %Thái Hòa-Nghệ An
\subsection{Trắc nghiệm (2 điểm)}
\Opensolutionfile{ans}[ans/ans8]
%Câu 1
\begin{ex}
Biểu thức $\dfrac{1}{\sqrt{3}+\sqrt{2}}+\dfrac{1}{\sqrt{3}-\sqrt{2}}$ có giá trị là:
\choice
{$1$}
{$3\sqrt{2}$}
{$\sqrt{3}$}
{\True $2\sqrt{3}$}
\end{ex}

%Câu 2
\begin{ex}
Cho 3 số thực $x,y,z$ biết $x<y$. Khẳng định nào sau đây \textbf{sai}?
\choice
{$x+z<y+z$}
{\True $xz<yz$ nếu $z<0$}
{$xz<yz$ nếu $z>0$}
{$x-z<y-z$}
\end{ex}

%Câu 3
\begin{ex}
Bất phương trình nào sau đây là bất phương trình bậc nhất một ẩn?
\choice
{$x-y<0$}
{$\sqrt{x}-5\le0$}
{$x^2-4\le0$}
{\True $\dfrac{5x-3}{2}>0$}
\end{ex}

%Câu 4
\begin{ex}
Cho tam giác $ABC$ vuông tại $A$, biết $AB=6$ cm, $AC=8$ cm. Khi đó $\sin\widehat{B}$ có giá trị là
\choice
{$\dfrac{4}{3}$}
{$\dfrac{3}{4}$}
{\True $\dfrac{4}{5}$}
{$\dfrac{6}{10}$}
\end{ex}

%Câu 5
\begin{ex}
Biểu thức $\sqrt{x-4}+\dfrac{\sqrt{4-x}}{\sqrt{x}}$ xác định khi
\choice
{\True $x=4$}
{$x>0$}
{$0<x\le 4$}
{$x\ge 4$}
\end{ex}

%Câu 6
\begin{ex}
Hệ thức nào sau đây sai?
\choice
{$\cos42^\circ=\sin48^\circ$}
{\True $\cot20^\circ=\tan80^\circ$}
{$\tan50^\circ=\cot40^\circ$}
{$\sin38^\circ=\cos52^\circ$}
\end{ex}

%Câu 7
\begin{ex}
Phản ứng giữa kim loại nhôm (Al) và oxygen ($\text{O}_{2}$) tạo thành hợp chất nhôm oxide ($\text{Al}_{2}\text{O}_{3}$) được biểu diễn bởi phương trình đã được cân bằng như sau:

\begin{center}
    $4\text{Al}+x\text{O}_{2} \rightarrow y\text{Al}_{2}\text{O}_{3}$
\end{center}
Khi đó $(x;y)$ bằng:
\choice
{$(2;3)$}
{\True $(3;2)$}
{$(3;3)$}
{$(2;2)$}
\end{ex}

%Câu 8
\begin{ex}
Cho đường tròn $(O;6\,\text{cm})$, đường kính $AB$. Gọi $M$ là một điểm trên đường tròn. Cho các khẳng định sau:
\begin{enumerate}[label=\roman*)]
\item Đường thẳng $AB$ là một trục đối xứng của đường tròn.
\item Tam giác $AMB$ là tam giác cân.
\item Nếu $\widehat{AOM}=60^\circ$ thì cung nhỏ $AM$ có độ dài bằng $2$ cm.
\item Nếu $MB=6\sqrt{3}$ cm thì $\triangle AOM$ đều.
\end{enumerate}
Hỏi có bao nhiêu khẳng định đúng?
\choice
{$1$}
{\True $2$}
{$0$}
{$3$}
\end{ex}

%Câu 9
\begin{ex}
Cho tam giác $ABC$ vuông tại $A$. Khẳng định nào sau đây đúng?
\choice
{$\tan B. \cos C=1$}
{\True $\sin^2 B + \cos^2 B=1$}
{$\sin A=\cos C$}
{$\tan^2 B+ \cot^2 C=1$}
\end{ex}

%Câu 10
\begin{ex}
Cho $\tan \alpha=2$. Tính $A=\dfrac{\sin \alpha +\cos \alpha}{\sin \alpha -\cos \alpha}$
\choice
{\True $3$}
{$2$}
{$1$}
{$-1$}
\end{ex}


\Closesolutionfile{ans}
\Opensolutionfile{ansbook}[ans/ansbook8]
\setcounter{bt}{0}
\subsection{Tự luận (8 điểm)}

%Bài 1
\begin{bt}[1,5 điểm] $    $
\begin{enumerate}
\item Giải phương trình: $\dfrac{1}{x+2}+\dfrac{4x}{x^2-4}=1+\dfrac{2}{x-2}$
\item Giải bất phương trình: $(x+2)^2-(3x-1)^2<-8x^2+3$
\end{enumerate}
\loigiai{
a) $x=1\,\text{và}\,x=2$ \hspace{2cm} b) $x<0$
}
\end{bt}

%Bài 2
\begin{bt}[1 điểm]
Hai đội công nhân cùng làm chung một công việc thì sau $12$ ngày sẽ xong. Nếu đội $1$ làm một mình trong $5$ ngày rồi nghỉ, đội $2$ làm tiếp trong $15$ ngày thì cả hai đội làm được $75\%$ công việc. Hỏi nếu làm một mình thì mỗi đội làm xong công việc đó trong bao lâu? 
\loigiai{Nếu làm một mình thì đội 1 hoàn thành công việc sau 20 ngày, đội 2 hoàn thành công việc sau 30 ngày.}
\end{bt}

%Bài 3
\begin{bt}[0,5 điểm]
\immini
{Một gia đình làm cầu thang có độ dốc là $36^\circ$ so với phương ngang. Chiều cao từ sàn nhà tới trần nhà là $3{,}6$ m. Tính chiều dài $AB$ của mặt cầu thang (làm tròn đến hàng đơn vị).}
{
\begin{tikzpicture}[scale=0.9, font=\small, >=stealth]
  \pgfmathsetmacro{\h}{3.6} % chiều cao
  \pgfmathsetmacro{\goc}{36} % góc nghiêng (độ)
  \pgfmathsetmacro{\day}{\h/tan(\goc)} % độ dài chân cầu thang
  \pgfmathsetmacro{\bac}{6} % số bậc
  \pgfmathsetmacro{\dx}{\day/\bac} % độ dài mỗi bậc
  \pgfmathsetmacro{\dy}{\h/\bac} % chiều cao mỗi bậc

  \coordinate (A) at (0,0);
  \coordinate (H) at (\day,0);
  \coordinate (B) at (\day,\h);

  \draw[thick] (A)--(B);
  \draw[dashed] (A)--(H) (B)--(H);

  % --- bậc thang: lưu ý có ; ở cuối \draw ---
  \foreach \i in {0,1,...,5}{
    \pgfmathsetmacro{\x}{\i*\dx}
    \pgfmathsetmacro{\y}{\i*\dy}
    \draw[thick] (\x,\y) -- ++(0,\dy) -- ++(\dx,0);
  }

  % góc ở A
  \pic[draw, angle radius=1cm, angle eccentricity=1.3] {angle = H--A--B};
  % nếu muốn nhãn rõ hơn, có thể dùng node riêng:
  % \node at ($(A)+(0.35,0.18)$) {$36^\circ$};

  % ký hiệu vuông góc tại H
  \draw ($(H)+(-0.3,0)$) -- ++(0,0.3) -- ++(0.3,0);

  % --- mũi tên chiều cao: phải vẽ thẳng đứng ---
  \draw[<->] ($(H)+(0.4,0)$) -- ++(0,\h) node[midway,right]{$3{,}6\ \mathrm{m}$};

  \foreach \p/\pos in {A/below,B/above right,H/below}{
    \filldraw[fill=black] (\p) circle(0.5pt) node[\pos] {$\p$};
  }
\end{tikzpicture}

}
\loigiai{
Chiều dài $AB$ của mặt cầu thang khoảng $6$ m.
}
\end{bt}

%Bài 4
\begin{bt}[1,5 điểm]
Rút gọn các biểu thức sau:
\begin{listEX}[2]
\item $A = 3\sqrt{3} + 2\sqrt{48} - \sqrt{75}$
\item $B = \dfrac{2}{\sqrt{3}-1}+\sqrt{(\sqrt{3}-4)^2}$
\item! $C = \left(\dfrac{\sqrt{x}}{3-\sqrt{x}} + \dfrac{x+9}{x-9}\right) : \dfrac{1}{3\sqrt{x}+x}$ \quad (với $x>0$, $x\ne9$)
\end{listEX}
\loigiai{
a) $A=6\sqrt{3}$ \hspace{2cm} b) $B=5$ \hspace{2cm} c) $C=-3\sqrt{x}$
}
\end{bt}

%Bài 5
\begin{bt} [3 điểm] Cho đường tròn $(O;R)$. Từ một điểm $M$ nằm ngoài đường tròn thỏa mãn $OM=2R$, vẽ các tiếp tuyến $MA$, $MB$ với đường tròn ($A,B$ là các tiếp điểm). Gọi $K$ là giao điểm của đoạn thẳng $OM$ và đường tròn $(O)$. 
\begin{enumerate}
    \item Tính số đo $\widehat{AOB}$ và $\widehat{AKB}$.
    \item Chứng minh rằng $AK$ là phân giác của $\widehat{MAB}$.
    \item Gọi $H$ là giao điểm của $AB$ và $OM$, $AK$ cắt $BM$ tại $I$. Chứng minh rằng $A,H,I,M$ cùng thuộc một đường tròn.
\end{enumerate}
 	\loigiai{
 		\immini{
 		\begin{enumerate}
 			\item $\widehat{AOB}=120^\circ,\widehat{AKB}=120^\circ$.
 			\item Chứng minh tam giác $MAB$ đều và $K$ là tâm đường tròn ngoại tiếp tam giác $MAB$.
 			\item Tam giác $MAB$ đều nên $AK$ vuông góc với $BM$. Hơn nữa $OM$ vuông góc với $AB$.
 		\end{enumerate}
 		}{
	\begin{tikzpicture}[line width=0.8pt, scale=1]
	\def\R{2.5}
	\def\A{5}
	\path 
	(0,0) coordinate (O)
	(\A,0) coordinate (M)
	($(O)!0.5!(M)$) coordinate (K);
	
	% Đường tròn tâm O và tâm K
	\path[name path=ciro] (O) circle (\R);
	\path[name path=cirom] (K) circle (\A/2); % bán kính KO = 2.5
	
	% Giao điểm của hai đường tròn
	\path[name intersections={of=ciro and cirom, by={A,B}}];
	
	% Vẽ hình
	\draw (O) circle (\R)
	      (A)--(O)--(B)
	      (O)--(M)
	      (M)--(B)--(A)--cycle;
	
	% Giao điểm, điểm phụ
	\path[name path=om] (O)--(M);
	\coordinate (H) at (intersection of A--B and O--M);
	\coordinate (At) at ($(A)!2!(K)$);
	\coordinate (I) at (intersection of A--At and B--M);
	\coordinate (Y) at ($(A)!0.5!(M)$);
	\draw (Y) let \p2=($(Y)-(A)$) in circle ({veclen(\x2,\y2)});
	\draw (A)--(K)--(B) (K)--(I)--(H);
	% Nhãn điểm
	\foreach \x/\g in {O/180,M/-40,A/90,B/-60,K/45,I/-90,H/-135}
	  \draw[fill=white] (\x) circle (1.5pt) +(\g:.3) node{$\x$};
\end{tikzpicture}
	}
  }
 \end{bt}

%Câu 6
\begin{bt}[0,5 điểm] Tính giá trị biểu thức $M=a^2+b^2$ biết $a,b$ thỏa mãn 
\[\begin{cases}
    \dfrac{3a^2}{b^2}+\dfrac{1}{b^3}=1\\
    \dfrac{3b^2}{a^2}+\dfrac{2}{a^3}=1
\end{cases} \, \text{với}\, a\ne 0, b\neq 0. \] 
\loigiai{
Biến đổi hệ phương trình ban đầu thành 
$\begin{cases}
    b^3-3a^2b=1\\
    a^3-3ab^2=2
\end{cases}$

Bình phương hai vế hai phương trình của hệ

Cộng hai vế của hai phương trình vừa tìm được ta được $(a^2+b^2)^3=5$

Từ đó suy ra $a^2+b^2=\sqrt[3]{5}$
}
\end{bt}
\Closesolutionfile{ansbook}